% 符号表段落模板（中文，自足性规则用——报告文首必含符号表）
% 用法：\input 在报告 \section{模型建立} 前；三列：符号/定义/首次出现处（节号或公式号）
% 自足性要求（TRAE-报告.md）：正文每个符号首次出现处一句话内嵌定义，此处汇总便于速查。

\begin{table}[ht]
\centering
\caption{模型符号说明（本文自足性符号表；正文首次出现处均有内嵌定义）}
\begin{tabular}{c l l}
\toprule
\textbf{符号} & \textbf{定义} & \textbf{首次出现} \\
\midrule
$\mathcal{R}$ & 六区域集合（RegionA--F） & §2.1 \\
$\mathcal{T}$ & 全时域 $\{0,\ldots,2405\}$（含结清段） & §2.1 \\
$A_j, D_j, g_j, L_j$ & 任务 $j$ 的到达小时/时长/GPU 需求/最晚完成 & §2.1 \\
% ……（每个符号一行；跨问代号如 S2 F4 也在此登记：代号|指代|首次出现）
\toprule
\textbf{跨问/实验代号} & \textbf{指代} & \textbf{首次出现} \\
\midrule
S2 F4 & S2 阶段 F4：全耦合变量核算 & §1.3 \\
Q3 对照 & 问题四 Q3：基荷策略合理性检验对照 & §3.2 \\
% ……（所有跨问引用/裁定代号逐一登记）
\bottomrule
\end{tabular}
\end{table}
